% ============================================================
%  Paper 5 — Aroha Labs
%  "The Aroha Protocol: A Layered Architecture for
%   Autonomous AI Agent Networks"
%
%  Target venues: AAMAS 2027 / WWW 2027 / IEEE Internet Computing
%  arXiv categories: cs.MA, cs.NI, cs.CR
% ============================================================

\documentclass[conference]{IEEEtran}
\IEEEoverridecommandlockouts

\usepackage{amsmath,amssymb,amsthm,mathtools}
\usepackage{algorithm}
\usepackage{algpseudocode}
\usepackage{graphicx}
\usepackage{booktabs}
\usepackage{hyperref}
\usepackage{xcolor}
\usepackage{cite}
\usepackage{bm}
\usepackage{listings}

\newtheorem{theorem}{Theorem}
\newtheorem{lemma}[theorem]{Lemma}
\newtheorem{proposition}[theorem]{Proposition}
\newtheorem{corollary}[theorem]{Corollary}
\newtheorem{definition}{Definition}
\newtheorem{remark}{Remark}

\DeclareMathOperator{\E}{\mathbb{E}}
\newcommand{\Adv}{\mathcal{A}}
\newcommand{\Sign}{\mathsf{Sign}}
\newcommand{\Verify}{\mathsf{Verify}}
\newcommand{\KeyGen}{\mathsf{KeyGen}}
\newcommand{\DID}{\mathsf{DID}}
\newcommand{\TL}{\mathsf{TL}}

\lstset{
  basicstyle={\ttfamily\scriptsize},
  breaklines=true,
  frame=none,
  columns=fullflexible
}

\begin{document}

\title{The Aroha Protocol: A Layered Architecture\\
       for Autonomous AI Agent Networks}

\author{
  \IEEEauthorblockN{Aroha Labs Research}
  \IEEEauthorblockA{%
    \textit{Aroha Protocol}\\
    aroha-labs.com
  }
}

\maketitle

% ─────────────────────────────────────────────────────────────
\begin{abstract}
The proliferation of large language model (LLM)-based agents
creates a demand for open, interoperable infrastructure that
enables agents published by different organisations to discover,
authenticate, communicate with, and delegate tasks to one another
with well-defined security guarantees. Existing agent frameworks
address single-provider deployments or rely on informal
conventions; no published work specifies a complete, layered
protocol for \emph{open} agent networks.
We present the \textit{Aroha Protocol}: a four-layer stack
comprising (L1) a content-addressed Decentralised Identifier
(DID) scheme derived from Ed25519 public keys, (L2) a dual
registry mechanism supporting both managed HTTP and on-chain
EVM-compatible registries, (L3) a signed capability manifest
system with IPFS-pinned content integrity, and (L4) a typed
envelope transport with a four-stage validation pipeline.
We formalise a trust-level taxonomy (TL0--TL3) that maps
cryptographic evidence to actionable trust decisions, prove the
security of key-rotation and key-revocation procedures, and
characterise the threat model for each layer. The protocol
includes bidirectional bridges to the Model Context Protocol
(MCP) and Agent-to-Agent (A2A) protocol, enabling incremental
adoption. A single signed envelope incurs approximately 650
bytes overhead; a three-agent delegation chain costs under 2 KB.
This paper is the architectural companion to four specialised
companion papers covering Bayesian reputation (P1), mandate
chains (P2), publisher verification (P3), and adaptive RAG
chunking (P4).
\end{abstract}

\begin{IEEEkeywords}
autonomous agents, decentralised identifiers, multi-agent
systems, protocol architecture, Ed25519, trust management,
agent interoperability, capability manifests
\end{IEEEkeywords}

% ─────────────────────────────────────────────────────────────
\section{Introduction}
\label{sec:intro}

Large language model (LLM)-based agents are rapidly transitioning
from isolated assistants to participants in collaborative,
multi-agent workflows~\cite{guo2024agents}. In such workflows an
agent may delegate sub-tasks to specialist agents, consume
published capabilities as services, or act as an orchestrator of
a pipeline spanning organisations and trust domains. This
transition imposes requirements that no existing framework
satisfies simultaneously: (i)~strong identity binding that
prevents impersonation, (ii)~capability advertisement that is
both machine-readable and tamper-evident, (iii)~a transport layer
whose security properties are independent of the underlying
network, and (iv)~interoperability with existing protocol
ecosystems such as MCP~\cite{anthropic2024mcp} and
A2A~\cite{google2024a2a}.

Existing work on multi-agent systems largely assumes a trusted
orchestrator or a closed environment. Reputation and trust models
for open peer-to-peer networks~\cite{josang2007trust,
kamvar2003eigentrust, sabater2005review} address agent quality
but not identity binding. Decentralised Identifiers
(DIDs)~\cite{w3c2022did} provide an identity layer but specify
no agent-specific capability model. Blockchain-based systems
offer tamper-resistant registries~\cite{nakamoto2008bitcoin,
buterin2014ethereum} but lack the application-level message
semantics required for agent coordination.

This paper presents the \textit{Aroha Protocol}: a complete,
layered specification for open AI agent networks. Our
contributions are:

\begin{enumerate}
  \item \textbf{A four-layer protocol stack} (DID scheme,
    registry, manifest, transport) with formal security
    properties at each layer (\S\ref{sec:stack}).
  \item \textbf{Theorems on key management security}: rotation
    cooldown prevents attacker exploitation; revocation
    self-signatures prove voluntary intent
    (\S\ref{sec:keymanagement}).
  \item \textbf{A trust-level taxonomy} (TL0--TL3) mapping
    cryptographic evidence to decision thresholds
    (\S\ref{sec:trust}).
  \item \textbf{A formal threat model} covering six attack
    classes and their mitigations (\S\ref{sec:security}).
  \item \textbf{Bidirectional bridges} to MCP and A2A with
    security downgrade characterisation
    (\S\ref{sec:bridges}).
\end{enumerate}

The protocol is designed so that each layer is independently
deployable; an implementor may adopt only the DID scheme and
transport layer while using an existing capability registry, or
adopt all four layers for a fully Aroha-native deployment.

% ─────────────────────────────────────────────────────────────
\section{Background and Related Work}
\label{sec:background}

\subsection{Decentralised Identifiers}

W3C DIDs~\cite{w3c2022did} decouple identifiers from
centralised registries by embedding verification material in a
DID Document. The \texttt{did:key} method derives a DID from a
public key, making it content-addressed and self-certifying.
Aroha adopts this principle but extends it with agent-specific
service endpoints, a structured capability manifest, and a
key-history field required for key rotation auditing.

\subsection{Trust and Reputation in Open Networks}

J\o{}sang and Ismail's Beta Reputation System~\cite{josang2002beta}
and EigenTrust~\cite{kamvar2003eigentrust} provide the
mathematical foundations for trust inference in open networks.
Our companion paper P1~\cite{aroha2025p1} extends these
foundations with asymmetric temporal decay tuned for LLM agents
whose quality evolves as underlying models are updated. The
present paper specifies the identity substrate on which the
reputation system operates.

\subsection{Cryptographic Delegation}

UCAN~\cite{fission2022ucan} and Macaroons~\cite{birgisson2014macaroons}
support capability delegation via chained tokens. Our companion
paper P2~\cite{aroha2025p2} defines mandate chains specifically
for financial delegation in agent workflows, building on the
Ed25519 primitive~\cite{bernstein2011ed25519,josefsson2017ed25519rfc}
specified in Aroha's identity layer.

\subsection{Content-Addressed Storage}

IPFS~\cite{ipfs2014whitepaper} provides content-addressed
immutable storage. Aroha uses IPFS CIDs (SHA-256 digests of
capability manifests) to bind registry references to manifest
content, preventing registry substitution attacks without
requiring on-chain storage of manifest data.

\subsection{Blockchain Registries}

Bitcoin~\cite{nakamoto2008bitcoin} and
Ethereum~\cite{buterin2014ethereum} demonstrate that distributed
ledger consensus can replace trusted intermediaries for state
agreement. Aroha's on-chain registry uses an EVM-compatible
smart contract, enabling deployment on Polygon PoS for
low-cost, high-throughput registration.

% ─────────────────────────────────────────────────────────────
\section{The Aroha Protocol Stack}
\label{sec:stack}

The Aroha Protocol is organised as four layers, each with a
defined interface to adjacent layers (Fig.~\ref{fig:stack}).
Layer dependencies flow upward: each layer assumes correct
operation of the layer below.

\begin{figure}[t]
\centering
\begin{tabular}{|c|}
\hline
\textbf{L4: Transport (ArohaEnvelope)} \\
\hline
\textbf{L3: Capability Manifest} \\
\hline
\textbf{L2: Registry} \\
\hline
\textbf{L1: DID Scheme} \\
\hline
\end{tabular}
\caption{The four-layer Aroha Protocol stack.}
\label{fig:stack}
\end{figure}

\subsection{Layer 1: DID Scheme}
\label{sec:l1}

\subsubsection{DID Syntax}

An Aroha DID takes one of three forms:

\begin{itemize}
  \item \textbf{Base form:}
    \texttt{did:aroha:\textit{<base58btc(pk)>}} where $pk$ is a
    32-byte Ed25519 public key encoded in Base58BTC~\cite{w3c2021multibase}.
  \item \textbf{Namespaced form:}
    \texttt{did:aroha:\textit{ns}.\textit{agentId}} for tenant
    isolation; the namespace portion is an organisation-controlled
    string.
  \item \textbf{Studio form:}
    \texttt{did:aroha:studio:\textit{h}} where $h =
    \mathrm{Keccak256}(\mathit{did\_string})$~\cite{keccak2015},
    used for agents registered via the Studio web interface.
\end{itemize}

\begin{definition}[Aroha DID]
Let $(\mathit{sk}, pk) \leftarrow \KeyGen()$ be an Ed25519
keypair. The canonical Aroha DID is:
\[
  \DID(pk) \;=\;
  \texttt{"did:aroha:"} \,\|\, \mathrm{Base58BTC}(pk).
\]
\end{definition}

Content-addressing ensures that $\DID(pk)$ is collision-free
under the hardness of discrete logarithm on Curve25519: two
distinct public keys produce two distinct DIDs.

\subsubsection{DID Document}

A DID Document for an Aroha agent contains:
(i)~\texttt{@context},
(ii)~\texttt{id} (the DID string),
(iii)~\texttt{verificationMethod} (type
\texttt{Ed25519VerificationKey2020}, \texttt{publicKeyMultibase}
prefixed with \texttt{z}),
(iv)~\texttt{authentication}, \texttt{assertionMethod},
\texttt{keyAgreement} (each referencing the verification method),
(v)~\texttt{service} (type \texttt{ArohaEndpoint}, pointing to
an HTTPS URL),
(vi)~\texttt{created}/\texttt{updated} timestamps, and
optionally (vii)~\texttt{keyHistory} and
\texttt{revokedKeys} arrays (see \S\ref{sec:keymanagement}).

\subsection{Layer 2: Registry}
\label{sec:l2}

\subsubsection{Registry Interface}

The canonical registry interface \texttt{IArohaRegistry}
provides:
\begin{itemize}
  \item \texttt{register(did, manifestCID, pk, opts)} —
    creates an agent record;
  \item \texttt{deactivate(did)} — marks the agent inactive;
  \item \texttt{updateManifest(did, newCID)} — updates the
    pinned manifest reference;
  \item \texttt{resolveByDID(did)} $\to$ \texttt{AgentRecord}
    $\mid$ null;
  \item \texttt{queryByCapability(ids[])} $\to$
    \texttt{AgentRecord[]};
  \item \texttt{onAgentRegistered(callback)} — event hook.
\end{itemize}

\subsubsection{HTTP Registry}

\texttt{ArohaHttpRegistry} exposes the interface as a managed
REST API. The registry is a trusted intermediary; clients must
verify manifest CIDs independently (see \S\ref{sec:l3}) because
HTTP transport does not provide consensus-based tamper resistance.

\subsubsection{On-Chain Registry}

\texttt{ArohaRegistryContract} is an EVM-compatible Solidity
smart contract deployed on Polygon PoS. Key design decisions:

\begin{itemize}
  \item \texttt{hashDID(did)} $\to$ \texttt{bytes32} and
    \texttt{hashCapability(id)} $\to$ \texttt{bytes32} map
    variable-length strings to fixed-width on-chain keys.
  \item A minimum stake requirement, enforced on-chain, raises
    the cost of Sybil attacks~\cite{douceur2002sybil}.
  \item The \texttt{AgentRegistered(didHash, owner,
    manifestCID)} event provides an append-only audit log.
\end{itemize}

Consensus-based tamper resistance means that modifying a
registry record requires 51\% of the PoS stake, which is
economically prohibitive for any single adversary.

\subsection{Layer 3: Capability Manifest}
\label{sec:l3}

\subsubsection{Manifest Structure}

A \textit{SignedCapabilityManifest} contains:
\begin{itemize}
  \item \texttt{agentDID} — the issuing agent's DID;
  \item \texttt{capabilities[]} — array of
    \texttt{CapabilityManifestEntry} records, each specifying a
    \texttt{name}, \texttt{description}, \texttt{inputSchema}
    (JSON Schema), optional \texttt{outputSchema},
    \texttt{supportsSaga} (bool), \texttt{defaultPricing}, and
    \texttt{defaultSLA};
  \item \texttt{generatedAt} — ISO 8601 timestamp;
  \item \texttt{signature} — Ed25519 signature over the
    canonical JSON~\cite{jsoncanon} of the manifest body.
\end{itemize}

The pricing model is one of \texttt{"per-call"},
\texttt{"subscription"}, or \texttt{"negotiated"}, with an
optional \texttt{amount} and \texttt{currency}. The SLA record
specifies \texttt{maxResponseMs} and an optional
\texttt{p99Ms} target.

\subsubsection{Content Integrity via IPFS CID}

The \textit{CID} (Content Identifier) of a manifest is:
\[
  \mathit{CID} = \mathrm{SHA}\text{-}256(\text{canonical JSON of manifest}).
\]
This CID is stored in the registry (both HTTP and on-chain)
and pinned to IPFS~\cite{ipfs2014whitepaper}. Any consumer
that resolves a registry record can re-derive the CID from the
fetched manifest and verify it matches the registry record;
a discrepancy proves a substitution attack.

\begin{theorem}[Manifest Integrity]
\label{thm:manifest}
Under SHA-256 collision resistance, no polynomial-time
adversary $\Adv$ can produce a manifest $M' \neq M$ such that
$\mathrm{SHA}\text{-}256(M') = \mathrm{SHA}\text{-}256(M)$.
Consequently, any manifest that passes CID verification is
identical to the manifest registered by the agent.
\end{theorem}
\begin{proof}
Immediate from the collision-resistance property of SHA-256
(standardised in FIPS 180-4 and confirmed by
NIST~\cite{fips186}).
\end{proof}

\subsection{Layer 4: Envelope Transport}
\label{sec:l4}

\subsubsection{ArohaEnvelope Structure}

Every inter-agent message is wrapped in an \textit{ArohaEnvelope}:
\begin{itemize}
  \item \texttt{@context}: \texttt{["https://aroha-labs.com/}\allowbreak\texttt{contexts/v1"]}
  \item \texttt{id}: UUID v4
  \item \texttt{type}: one of 16 \texttt{ArohaMessageType} values
  \item \texttt{from}, \texttt{to}: DID strings
  \item \texttt{created}, \texttt{expires}: ISO 8601 timestamps
  \item \texttt{nonce}: 128-bit random value (hex string)
  \item \texttt{correlationId}: links messages in a saga
  \item \texttt{body}: typed payload (per message type)
  \item \texttt{proof?}: \texttt{MessageProof} (Ed25519
    signature over canonical JSON of envelope minus
    \texttt{proof} field)
  \item \texttt{traceparent?}: W3C Trace Context Level 1
    header~\cite{w3c2021tracecontext}
  \item \texttt{routingHint?}: complexity indicator and optional
    model suggestion
\end{itemize}

\subsubsection{Envelope Size}

Fixed overhead fields have bounded sizes: an Ed25519 public key
in Base58BTC occupies approximately 44 characters; a signature
approximately 88 characters; a nonce 22 characters; a UUID 36
characters. A typical \texttt{ArohaRequest} envelope with a
small body occupies approximately 650 bytes. A three-agent
delegation chain (orchestrator $\to$ worker $\to$ sub-worker)
produces three envelopes totalling approximately 1.95 KB,
which is negligible relative to LLM prompt or response payloads.

\subsubsection{Envelope Validation Pipeline}

Algorithm~\ref{alg:validate} specifies the four-stage sequential
validation procedure executed by every receiving agent.

\begin{algorithm}[t]
\caption{Envelope Validation Pipeline}
\label{alg:validate}
\begin{algorithmic}[1]
\Require Envelope $E$, local DID $\mathit{myDID}$,
         nonce registry $\mathcal{N}$,
         options $\mathit{opts}$
\Ensure \textsc{Accept} or \textsc{Reject}
\State \textbf{Stage 1: Signature Verification}
\If{$\mathit{opts.bypassSignatureFor}(E.\texttt{from})$}
  \State proceed \Comment{Trusted Mesh bypass}
\Else
  \State $pk \leftarrow \mathrm{resolveKey}(E.\texttt{from})$
  \State $m \leftarrow \mathrm{canonicalJSON}(E \setminus
         \{\texttt{proof}\})$
  \If{$\lnot\,\Verify(pk,\, m,\, E.\texttt{proof.sig})$}
    \State \Return \textsc{Reject}(\textit{bad-sig})
  \EndIf
\EndIf
\State \textbf{Stage 2: Expiry Check}
\State $\delta \leftarrow \mathit{opts.clockToleranceMs}$
       $\!\parallel\! 0$
\If{$E.\texttt{expires} + \delta \leq \mathrm{now()}$}
  \State \Return \textsc{Reject}(\textit{expired})
\EndIf
\State \textbf{Stage 3: Nonce Freshness}
\If{$E.\texttt{nonce} \in \mathcal{N}$}
  \State \Return \textsc{Reject}(\textit{replay})
\EndIf
\State $\mathcal{N}.\mathrm{insert}(E.\texttt{nonce})$
\State \textbf{Stage 4: Recipient Match}
\If{$E.\texttt{to} \neq \mathit{myDID}$}
  \State \Return \textsc{Reject}(\textit{misdirected})
\EndIf
\State \Return \textsc{Accept}
\end{algorithmic}
\end{algorithm}

The four stages are independent and ordered by computational
cost: recipient-match (O(1) string comparison) is placed last
because it catches only routing errors, while signature
verification (the most expensive stage) is placed first to
avoid committing nonce state for forged envelopes.

% ─────────────────────────────────────────────────────────────
\section{Key Management}
\label{sec:keymanagement}

\subsection{Key Rotation}

Key rotation is triggered by the owner via
$\mathrm{rotateDIDKey}(\mathit{current}, \Delta)$, where
$\Delta$ is a cooldown duration in hours. The procedure generates
a new keypair $(\mathit{sk}', pk')$, creates a
\textit{KeyRotationRecord}:
{\small\begin{align*}
\bigl\{\;
  &\mathit{rotatedAt},\;
  \mathit{trustAfter} = \mathit{rotatedAt} + \Delta,\;\\
  &\mathit{predecessorKey},\;
  \mathit{newKeyHash},\;
  \sigma_{\mathit{pred}}
\;\bigr\}
\end{align*}}
where $\mathit{newKeyHash} = \mathrm{SHA}\text{-}256(pk')$ and
the predecessor signature is:
{\small\begin{align*}
  \sigma_{\mathit{pred}} = \Sign\bigl(\mathit{sk},\;
    \mathit{predecessorKey} \,\|\,
    \mathit{newKeyHash} \,\|\,
    \mathit{rotatedAt}\bigr).
\end{align*}}
The function $\mathrm{extractPublicKey}(\mathit{doc})$ raises an
exception if $\mathrm{now}() < \mathit{trustAfter}$.

\begin{theorem}[Rotation Cooldown Security]
\label{thm:rotation}
Let $\Adv$ be a computationally bounded adversary who obtains
temporary read access to $\mathit{sk}$ during interval
$[t_0, t_1]$ and initiates key rotation at time $t_r \in
[t_0, t_1]$ with cooldown $\Delta > 0$. Then $\Adv$'s new key
$pk'_{\Adv}$ is untrusted until $t_r + \Delta$, during which
time the legitimate owner can detect the unauthorised rotation
(via the $\mathit{keyHistory}$ audit log) and issue an
emergency revocation. Specifically, if the owner checks
$\mathit{keyHistory}$ at any time $t \in (t_r, t_r + \Delta)$,
the adversary's key is never trusted.
\end{theorem}
\begin{proof}
By the protocol, $\mathrm{extractPublicKey}$ enforces
$\mathrm{now}() \geq \mathit{trustAfter} = t_r + \Delta$ before
returning the new key. Since the predecessor signature
$\sigma_{\mathit{pred}}$ is included in the rotation record and
the \texttt{keyHistory} field is append-only (enforced by the
registry), the legitimate owner observing any key in
$\mathit{keyHistory}$ can identify rotation records initiated
outside their control. Issuing $\mathrm{revokeKey}$ before
$t_r + \Delta$ permanently invalidates $pk'_{\Adv}$ without a
cooldown (revocation is immediate by design). Therefore, provided
the owner monitors $\mathit{keyHistory}$ at frequency
$f > 1/\Delta$, the adversary's window of trust is the empty set.
\end{proof}

\subsection{Key Revocation}

Emergency revocation is performed by
$\mathrm{revokeKey}(\mathit{current}, \mathit{reason})$, which
creates a \textit{RevokedKeyRecord}:
\[
\{
  \mathit{keyMultibase},\;
  \mathit{revokedAt},\;
  \mathit{reason},\;
  \sigma_{\mathit{self}}
\}
\]
where the self-signature is:
{\small\begin{align*}
  \sigma_{\mathit{self}} = \Sign\bigl(\mathit{sk}_{\mathit{rev}},\;
    \mathit{keyMultibase} \,\|\,
    \mathit{revokedAt} \,\|\,
    \mathit{reason}\bigr).
\end{align*}}

\begin{theorem}[Voluntary Revocation Proof]
\label{thm:revocation}
Under the unforgeability of Ed25519 signatures, a valid
$\sigma_{\mathit{self}}$ over a \textit{RevokedKeyRecord}
constitutes a proof that the holder of $\mathit{sk}_{\mathit{rev}}$
voluntarily issued the revocation. No party lacking
$\mathit{sk}_{\mathit{rev}}$ can produce a valid
$\sigma_{\mathit{self}}$ with non-negligible probability.
\end{theorem}
\begin{proof}
Follows directly from the existential unforgeability under
chosen-message attack (EUF-CMA) security of Ed25519 as proved
in~\cite{bernstein2011ed25519}. An adversary who does not
possess $\mathit{sk}_{\mathit{rev}}$ cannot produce a valid
signature on the concatenation
$\mathit{keyMultibase} \,\|\, \mathit{revokedAt} \,\|\,
\mathit{reason}$ without breaking the discrete logarithm
assumption on Curve25519.
\end{proof}

The self-signature serves a practical purpose: it distinguishes
a voluntary revocation (issued by the key's owner) from an
attacker-forced revocation, which the attacker could mount only
if they already possess the private key. In the latter case the
attacker's incentive is to \emph{use} the key rather than revoke
it, so a valid self-signature provides meaningful assurance in
practice.

% ─────────────────────────────────────────────────────────────
\section{Trust Level Taxonomy}
\label{sec:trust}

We formalise a four-level trust taxonomy that maps the
cryptographic evidence attached to an incoming message to a
decision threshold. Table~\ref{tab:trustlevels} summarises the
levels.

\begin{table}[t]
\centering
\caption{Aroha Trust Level Taxonomy}
\label{tab:trustlevels}
\begin{tabular}{@{}clll@{}}
\toprule
\textbf{Level} & \textbf{DID} & \textbf{Signature} &
  \textbf{Typical Source} \\
\midrule
TL0 & None & None & Unauthenticated, A2A bridge \\
TL1 & Registered & None & Trusted Mesh peers \\
TL2 & Registered & Verified & Standard Aroha agents \\
TL3 & Registered & Verified & Verified-publisher (P3) \\
\bottomrule
\end{tabular}
\end{table}

\begin{definition}[Trust Level]
Given an incoming envelope $E$, define:
\begin{align*}
\TL(E) = \begin{cases}
  3 & \text{if } \TL(E) \geq 2
      \text{ and publisher-verified (P3)} \\
  2 & \text{if DID registered} \land \Verify(\cdot) = \top \\
  1 & \text{if DID registered} \land \text{no sig or bypassed} \\
  0 & \text{otherwise}
\end{cases}
\end{align*}
\end{definition}

\begin{remark}
TL1 applies within Trusted Mesh mode (\S\ref{sec:mesh}), where
network-level controls (VPC or mTLS service mesh) substitute for
per-message signatures. Nonce and expiry checks remain enforced
at TL1; only the signature check is bypassed for explicitly
allowlisted DIDs.
\end{remark}

Downstream policy decisions (e.g., whether to honour a financial
mandate, whether to execute a destructive capability) are
conditioned on $\TL(E)$ meeting a caller-configured threshold.
The companion paper P2~\cite{aroha2025p2} requires $\TL \geq 2$
for mandate chain execution; the publisher verification system
in P3~\cite{aroha2025p3} grants elevated search ranking only to
$\TL = 3$ agents.

% ─────────────────────────────────────────────────────────────
\section{Trusted Mesh Mode}
\label{sec:mesh}

Within a homogeneous deployment (e.g., all agents in a single
Kubernetes namespace or VPC), per-message Ed25519 signatures
impose redundant overhead because the transport layer already
provides mutual authentication (mTLS) and integrity (TLS
record-layer MAC). Trusted Mesh mode eliminates this overhead
for intra-mesh communication.

Calling \texttt{createTrustedMeshOptions} with a set of
trusted DIDs and optional namespace prefixes returns a
validation options object
with:
{\small\begin{align*}
  \mathit{bypassSig}(\mathit{did}) &=
    \mathit{did} \in \mathit{trustedDIDs} \\
    &\;\lor\;
    \exists\, p \in \mathit{trustedPrefixes}:\,
      \mathit{did.startsWith}(p).
\end{align*}}

The \textit{TrustedMeshClient} sends unsigned envelopes via
$\mathrm{buildUnsignedEnvelope}()$; receiving agents skip Stage~1
of Algorithm~\ref{alg:validate} for bypassed DIDs while
enforcing Stages~2--4. External agents (not in
$\mathit{trustedDIDs}$) cannot bypass signature verification
even if they can reach the internal endpoint.

$\mathrm{createTrustedZone}(\mathit{trustedDIDs},
\mathit{prefixes?})$ packages server-side validation options and
a pre-configured client factory into a named zone object,
simplifying mesh topology management.

% ─────────────────────────────────────────────────────────────
\section{Threat Model and Security Analysis}
\label{sec:security}

Table~\ref{tab:threats} enumerates the six principal attack
classes against the Aroha Protocol, the layer at which each is
mitigated, and the mechanism employed.

\begin{table}[t]
\centering
\caption{Aroha Threat Model}
\label{tab:threats}
\renewcommand{\arraystretch}{1.1}
\begin{tabular}{@{}p{2.1cm}p{1.1cm}p{4.5cm}@{}}
\toprule
\textbf{Attack} & \textbf{Layer} & \textbf{Mitigation} \\
\midrule
DID spoofing &
  L1 &
  DID is derived from public key; forging a DID requires
  finding a key pair with the target public key (DL hardness)
  \\[2pt]
Key compromise &
  L1 &
  Rotation cooldown + revocation; adversary's new key is
  untrusted for $\Delta$ hours (Theorem~\ref{thm:rotation})
  \\[2pt]
Registry substitution (HTTP) &
  L3 &
  Manifest CID re-derived at consumer; mismatch detected
  (Theorem~\ref{thm:manifest})
  \\[2pt]
Registry tampering (chain) &
  L2 &
  51\% PoS consensus required; economically prohibitive
  \\[2pt]
Replay attack &
  L4 &
  Nonce registry rejects duplicate nonces within rolling
  window; \texttt{expires} provides absolute bound
  \\[2pt]
Routing / impersonation &
  L4 &
  \texttt{to} field check (Stage~4) rejects envelopes
  addressed to another DID
  \\
\bottomrule
\end{tabular}
\end{table}

\subsection{Sybil Resistance}

The on-chain registry's staking requirement raises the cost
of Sybil attacks~\cite{douceur2002sybil} linearly with the
number of registered identities. The HTTP registry relies on
API authentication and rate limiting; operators may impose
identity verification requirements at registration time, which
is the mechanism formalised in companion paper
P3~\cite{aroha2025p3}.

\subsection{Canonical JSON and Signature Malleability}

Signatures are computed over canonical JSON~\cite{jsoncanon}
(RFC~8785 JCS) to eliminate signature malleability attacks that
exploit JSON key-ordering ambiguity. All fields in the envelope
except \texttt{proof} are included; the \texttt{proof} field is
excluded to avoid circular dependencies.

% ─────────────────────────────────────────────────────────────
\section{Protocol Bridges}
\label{sec:bridges}

The Aroha Protocol does not require all participants to be
Aroha-native. Two bidirectional bridges enable interoperability
with the dominant existing ecosystems.

\subsection{MCP Bridge}

The Model Context Protocol~\cite{anthropic2024mcp} defines a
tool-call interface between LLM hosts and tool servers.
Originally introduced by Anthropic, MCP was donated to the
Linux Foundation's Agentic AI Foundation in December 2025 and
has since become the de~facto industry standard for
LLM--tool integration, making bridge fidelity a first-order
interoperability requirement rather than a convenience.
The Aroha--MCP bridge operates in both directions:
\begin{itemize}
  \item \textbf{Aroha $\to$ MCP:}
    $\mathrm{arohaAgentToMcpTools}()$ wraps each entry in an
    agent's capability manifest as an MCP tool definition.
    The \texttt{inputSchema} field maps directly to the MCP
    tool's JSON Schema parameter specification.
  \item \textbf{MCP $\to$ Aroha:}
    $\mathrm{mcpToolsToArohaCapabilities}()$ exposes an MCP
    server as a registered Aroha agent. MCP tool calls are
    translated to \texttt{ArohaRequest} envelopes and executed
    against the MCP server.
\end{itemize}

Because MCP tool calls carry full request semantics and the
Aroha envelope wraps the call end-to-end, the MCP bridge
preserves the TL2 trust level of the originating agent.

\subsection{A2A Bridge}

Google's Agent-to-Agent protocol~\cite{google2024a2a} specifies
an HTTP-based message-passing protocol for LLM agents with its
own task-tracking semantics. The Aroha--A2A bridge translates:
\begin{itemize}
  \item A2A \texttt{messageId} $\leftrightarrow$ Aroha
    \texttt{id} (UUID v4);
  \item A2A \texttt{taskId} $\leftrightarrow$ Aroha
    \texttt{correlationId};
  \item A2A \texttt{role=user} $\leftrightarrow$ Aroha
    \texttt{ArohaRequest};
  \item A2A \texttt{role=agent} $\leftrightarrow$ Aroha
    \texttt{ArohaResponse}.
\end{itemize}

\textbf{Security downgrade:} A2A does not carry Ed25519
signatures. Inbound A2A messages therefore receive trust level
TL1 (registered DID, signature absent) rather than the TL2
they would receive as native Aroha messages. Implementations
MUST NOT execute capabilities that require $\TL \geq 2$ for
messages originating from the A2A bridge unless an out-of-band
trust agreement exists.

\subsection{Additional Bridges}

The Aroha SDK includes bridges for LangChain~\cite{harrison2022langchain},
Hermes, and Composio. These follow the same pattern: map the
external framework's tool or action abstraction to Aroha
capability entries, and wrap invocations in
\texttt{ArohaRequest} envelopes. Trust levels mirror A2A
(TL0--TL1) unless the external framework provides its own
signature mechanism.

Table~\ref{tab:bridges} summarises the security properties of
each bridge.

\begin{table}[t]
\centering
\caption{Bridge Security Properties}
\label{tab:bridges}
\begin{tabular}{@{}lccl@{}}
\toprule
\textbf{Bridge} & \textbf{Bidirectional} &
  \textbf{Min. TL} & \textbf{Notes} \\
\midrule
MCP           & Yes & TL2 & Preserves Ed25519 chain \\
A2A           & Yes & TL1 & No sig; DID-registered only \\
LangChain     & Yes & TL0 & Framework-level auth only \\
Hermes        & Yes & TL0 & Framework-level auth only \\
Composio      & Yes & TL0 & Framework-level auth only \\
\bottomrule
\end{tabular}
\end{table}

% ─────────────────────────────────────────────────────────────
\section{Integration with Companion Papers}
\label{sec:companion}

The Aroha Protocol stack is designed so that each companion
paper adds a vertical cross-cutting concern rather than a new
layer:

\begin{itemize}
  \item \textbf{P1 — ArohaRep~\cite{aroha2025p1}:} The
    reputation engine operates on agent DIDs resolved via L1
    and L2. Manifest integrity (Theorem~\ref{thm:manifest})
    provides signal $s_{10}$ in ArohaRep's eleven-signal
    observation model. Discovery queries are filtered through
    ArohaRep's Thompson Sampling ranking.
  \item \textbf{P2 — ArohaMandate~\cite{aroha2025p2}:} Mandate
    chains are encoded in \texttt{ArohaRequest} envelope bodies
    and validated at L4 Stage~1. The $\TL \geq 2$ requirement
    for mandate execution is enforced by the validation pipeline
    (Algorithm~\ref{alg:validate}).
  \item \textbf{P3 — Publisher Verification~\cite{aroha2025p3}:}
    The tiered publisher verification system issues a
    \textit{publisher credential} that elevates an agent's
    trust level to TL3. The credential is stored in the agent's
    DID Document service endpoint and verified at L4
    Stage~1.
  \item \textbf{P4 — ArohaChunk~\cite{aroha2025p4}:} The
    adaptive RAG chunking system~\cite{lewis2020rag} operates
    on document corpora retrieved via \texttt{ArohaRequest}
    envelopes. Chunk boundaries and retrieval results are
    returned as \texttt{ArohaResponse} bodies, subject to the
    same validation pipeline.
\end{itemize}

% ─────────────────────────────────────────────────────────────
\section{Implementation Notes}
\label{sec:impl}

The reference implementation is a TypeScript/Node.js SDK. Key
implementation decisions include:

\subsubsection{Nonce Registry}
The \texttt{NonceRegistry} uses a time-bucketed rolling window
keyed by nonce (128-bit random hex). Expired buckets are pruned
lazily on each insertion. For clustered deployments, a shared
Redis instance provides a distributed nonce store with TTL-based
expiry.

\subsubsection{Canonical JSON}
The SDK uses RFC~8785 JCS~\cite{jsoncanon} via a pure-TypeScript
implementation. The canonicalisation is applied once per
outbound message and once per inbound signature verification;
the result is not cached to avoid state-dependent bugs.

\subsubsection{Registry Multiplexing}
An \texttt{ArohaCompositeRegistry} wraps both HTTP and on-chain
registries behind the \texttt{IArohaRegistry} interface, routing
writes to the on-chain registry and reads to whichever registry
responds first, with CID cross-verification.

\subsubsection{Multibase Encoding}
Public keys and revocation records use the
\texttt{z}-prefixed Base58BTC multibase
encoding~\cite{w3c2021multibase}, consistent with the W3C DID
Core recommendation~\cite{w3c2022did}. The \texttt{z} prefix
unambiguously identifies the encoding scheme, enabling
future algorithm agility.

% ─────────────────────────────────────────────────────────────
\section{Discussion}
\label{sec:discussion}

\subsection{Comparison with Existing Agent Protocols}

Unlike MCP~\cite{anthropic2024mcp} and A2A~\cite{google2024a2a},
the Aroha Protocol specifies a complete identity layer (L1) and
a capability registry (L2--L3) in addition to a message transport
(L4). MCP intentionally omits identity; A2A includes task state
but not cryptographic identity binding. Aroha is complementary:
it provides the identity and security substrate on which either
protocol can operate via the bridges in \S\ref{sec:bridges}.

\subsection{Limitations}

\paragraph{Key monitoring obligation.}
Theorem~\ref{thm:rotation} requires the legitimate owner to
monitor \texttt{keyHistory} at frequency $f > 1/\Delta$. An
owner who does not monitor may miss an adversarial rotation
within the cooldown window. The SDK provides a
\texttt{KeyRotationMonitor} helper that polls the registry and
raises an alert for unexpected rotation records; this shifts the
obligation from passive inspection to active monitoring.

\paragraph{On-chain latency.}
The Polygon PoS network imposes a confirmation latency of
approximately 2--5 seconds per block. Applications requiring
sub-second registration confirmation must use the HTTP registry
and accept its weaker tamper-resistance guarantees.

\paragraph{Trusted Mesh boundary enforcement.}
The Trusted Mesh mode relies on the network operator correctly
configuring access controls so that external agents cannot reach
internal endpoints. A misconfigured mesh (e.g., a firewall rule
that inadvertently exposes the internal endpoint to the internet)
would allow external agents to send unsigned envelopes that
receive TL1 treatment.

\subsection{Future Work}

Several directions extend the Aroha Protocol:
(i)~post-quantum signature algorithms (CRYSTALS-Dilithium)
as a drop-in replacement for Ed25519, requiring only a
new \texttt{verificationMethod} type;
(ii)~zero-knowledge proofs of capability possession to enable
privacy-preserving discovery;
(iii)~threshold signatures for multi-party agent governance,
where a capability manifest requires $k$-of-$n$ organisational
signatures before activation.

% ─────────────────────────────────────────────────────────────
\section{Conclusion}
\label{sec:conclusion}

We have presented the Aroha Protocol: a four-layer architecture
for open AI agent networks comprising a content-addressed DID
scheme, dual HTTP/on-chain registry implementations, a
IPFS-pinned capability manifest system, and a typed envelope
transport with a four-stage validation pipeline. We proved that
key rotation with a cooldown parameter mitigates temporary key
compromise (Theorem~\ref{thm:rotation}), that Ed25519
self-signatures on revocation records provide voluntary-intent
proofs (Theorem~\ref{thm:revocation}), and that SHA-256 CID
verification provides manifest integrity against substitution
attacks (Theorem~\ref{thm:manifest}). We formalised a
four-level trust taxonomy (TL0--TL3) and characterised the
security properties of bidirectional bridges to MCP and A2A.

The Aroha Protocol addresses a gap in the current ecosystem: no
prior specification provides a complete, layered, cryptographically
grounded identity and communication substrate for open AI agent
networks. We believe the protocol—and its companion papers
covering reputation (P1), financial delegation (P2), publisher
verification (P3), and document intelligence (P4)—form a
coherent, deployable foundation for the next generation of
autonomous agent infrastructure.

% ─────────────────────────────────────────────────────────────
\bibliographystyle{IEEEtran}
\bibliography{aroha}

\end{document}
